% Exemple: dos límits calculats amb la definició (valencià).

\begin{frame}
\didactatitle{Un límit per la definició}

\begin{example}
Provem que $\lim_{x \to 2} (3x - 1) = 5$. Donat $\varepsilon > 0$, busquem
$\delta > 0$ tal que $0 < |x - 2| < \delta$ implique
$|(3x - 1) - 5| < \varepsilon$. Com que
\[
  |(3x - 1) - 5| = |3x - 6| = 3\,|x - 2| ,
\]
n'hi ha prou de prendre $\delta = \varepsilon / 3$: si
$0 < |x - 2| < \varepsilon/3$, aleshores $|(3x - 1) - 5| = 3\,|x - 2| <
\varepsilon$.
\end{example}

\onlynotes{%
  El mètode és sempre el mateix: escriure $|f(x) - L|$ en termes de
  $|x - a|$ i llegir-hi quant cal acostar-se.
}
\end{frame}

\begin{frame}
\didactatitle{Fitar abans de triar}

\begin{example}
Provem que $\lim_{x \to 1} x^2 = 1$. Ara
\[
  |x^2 - 1| = |x - 1|\,|x + 1| ,
\]
i el factor $|x + 1|$ depén de $x$. El fitem: si $|x - 1| < 1$, aleshores
$0 < x < 2$ i $|x + 1| < 3$, de manera que $|x^2 - 1| < 3\,|x - 1|$.

\dpause

Prenem
\begin{keyformula}
  \delta = \min\{1, \varepsilon/3\} .
\end{keyformula}
Si $0 < |x - 1| < \delta$, es compleixen alhora $|x - 1| < 1$ i
$|x - 1| < \varepsilon/3$, per tant $|x^2 - 1| < 3\,|x - 1| < \varepsilon$.
\end{example}

\begin{commonmistake}
Prenen $\delta = \varepsilon / |x + 1|$. Però $\delta$ es tria abans que $x$:
pot dependre de $\varepsilon$ i de $a$, mai de $x$.
\end{commonmistake}
\end{frame}
